\chapter{Linear Mappings}

\section{Matrix mappings}

\begin{definition}[Matrix mapping]
For an $m\times n$ matrix $A$, the function $f_A:\R^n\to\R^m$ defined by $f_A(\vec x)=A\vec x$ is the
\emph{matrix mapping} corresponding to $A$. $\R^n$ is the \emph{domain} and $\R^m$ the
\emph{codomain}.
\end{definition}

By linearity of matrix--vector multiplication, every matrix mapping satisfies
$f_A(b\vec x+c\vec y)=bf_A(\vec x)+cf_A(\vec y)$. This property is what we abstract.

\begin{definition}[Linear mapping]
A function $L:\R^n\to\R^m$ is a \emph{linear mapping} if for all $\vec x,\vec y\in\R^n$ and
$b,c\in\R$,
\[
  L(b\vec x + c\vec y) = bL(\vec x)+cL(\vec y).
\]
A linear mapping $L:\R^n\to\R^n$ is called a \emph{linear operator}.
\end{definition}

Immediate consequences: $L(\vec 0)=\vec 0$ (take $b=c=0$), and
$L(c_1\vec v_1+\dots+c_k\vec v_k)=c_1L(\vec v_1)+\dots+c_kL(\vec v_k)$.

\begin{example}
$L(x_1,x_2,x_3)=(x_1+2x_2,\ 3x_2-x_3)$ is linear: for $\vec x,\vec y$ and $b,c$,
\[
L(b\vec x+c\vec y)=\big(bx_1+cy_1+2bx_2+2cy_2,\ 3bx_2+3cy_2-bx_3-cy_3\big)=bL(\vec x)+cL(\vec y).
\]
In contrast, $M(x_1,x_2)=(x_1+1,x_2)$ is not linear since $M(\vec0)=(1,0)\neq\vec 0$, and
$N(x_1,x_2)=(x_1x_2, 0)$ is not linear since $N(2(1,1))=(4,0)\ne 2N(1,1)=(2,0)$.
\end{example}

\section{The standard matrix}

\begin{theorem}
Every linear mapping $L:\R^n\to\R^m$ is a matrix mapping: $L(\vec x)=[L]\vec x$ where
\[
  [L] = \big[L(\vec e_1)\ \ L(\vec e_2)\ \cdots\ L(\vec e_n)\big]
\]
is the \emph{standard matrix} of $L$. It is the unique matrix with this property.
\end{theorem}
\begin{proof}
$L(\vec x)=L(x_1\vec e_1+\dots+x_n\vec e_n)=x_1L(\vec e_1)+\dots+x_nL(\vec e_n)=[L]\vec x$ by the
definition of matrix--vector multiplication. Uniqueness is the Matrices Equal Theorem.
\end{proof}

\begin{example}
For $L(x_1,x_2,x_3)=(x_1+2x_2,\ 3x_2-x_3)$: $L(\vec e_1)=(1,0)$, $L(\vec e_2)=(2,3)$,
$L(\vec e_3)=(0,-1)$, so
$[L]=\begin{bmatrix}1&2&0\\0&3&-1\end{bmatrix}$.
\end{example}

So ``linear mapping $\R^n\to\R^m$'' and ``$m\times n$ matrix'' are two views of the same object.

\section{Geometric transformations}

\subsection*{Rotations}
$R_\theta:\R^2\to\R^2$ rotates counter-clockwise by $\theta$. Rotation sends $\vec e_1$ to
$(\cos\theta,\sin\theta)$ and $\vec e_2$ to $(-\sin\theta,\cos\theta)$, and (geometrically) preserves
sums and scalings, so it is linear with
\[
  [R_\theta]=\begin{bmatrix}\cos\theta&-\sin\theta\\ \sin\theta&\cos\theta\end{bmatrix}.
\]
For example $[R_{\pi/2}]=\begin{bmatrix}0&-1\\1&0\end{bmatrix}$.

\subsection*{Projections}
For $\vec d\ne\vec 0$, $\proj_{\vec d}$ is linear (Chapter 2). For $\vec d=(1,2)$,
$\proj_{\vec d}\vec e_1=\frac15(1,2)$ and $\proj_{\vec d}\vec e_2=\frac25(1,2)$, so
\[
  [\proj_{\vec d}]=\frac15\begin{bmatrix}1&2\\2&4\end{bmatrix}.
\]
In general $[\proj_{\vec d}]=\dfrac{1}{\norm{\vec d}^2}\vec d\vec d^{\,T}$.

\subsection*{Reflections}
Reflecting $\vec x$ in the line $\Span\{\vec d\}$ keeps the projection and flips the perpendicular part:
$\refl_{\vec d}\vec x = \proj_{\vec d}\vec x-\perpp_{\vec d}\vec x = 2\proj_{\vec d}\vec x-\vec x$. For
$\vec d=(1,2)$,
\[
  [\refl_{\vec d}] = \frac25\begin{bmatrix}1&2\\2&4\end{bmatrix}-I = \frac15\begin{bmatrix}-3&4\\4&3\end{bmatrix}.
\]
Check: $\frac15(-3+8,\ 4+6)=(1,2)$, so $\vec d$ is fixed, and $(2,-1)\mapsto\frac15(-6-4,\ 8-3)=(-2,1)$
is flipped. A reflection in a plane through the origin in $\R^3$ with normal $\vec n$ is
$\refl_P\vec x = \vec x - 2\proj_{\vec n}\vec x$.

\subsection*{Stretches and shears}
$\begin{bmatrix}t&0\\0&1\end{bmatrix}$ stretches horizontally by $t$; $\begin{bmatrix}1&s\\0&1\end{bmatrix}$
is a horizontal shear, $(x_1,x_2)\mapsto(x_1+sx_2,x_2)$.

\section{Operations on linear mappings and composition}

\begin{definition}
For $L,M:\R^n\to\R^m$ and $c\in\R$, define $(L+M)(\vec x)=L(\vec x)+M(\vec x)$ and
$(cL)(\vec x)=cL(\vec x)$. For $M:\R^n\to\R^m$ and $L:\R^m\to\R^p$, the \emph{composition} is
$(L\circ M)(\vec x)=L(M(\vec x))$.
\end{definition}

\begin{theorem}
These are linear, and $[L+M]=[L]+[M]$, $[cL]=c[L]$, $[L\circ M]=[L][M]$.
\end{theorem}
\begin{proof}[Proof for composition]
$(L\circ M)(\vec x)=[L]([M]\vec x)=([L][M])\vec x$ by associativity, so $L\circ M$ is the matrix
mapping of $[L][M]$, which is linear.
\end{proof}

\begin{example}
Rotate by $\pi/2$ and then reflect in the $x_1$-axis ($F=\begin{bmatrix}1&0\\0&-1\end{bmatrix}$):
\[
  [F\circ R_{\pi/2}] = \begin{bmatrix}1&0\\0&-1\end{bmatrix}\begin{bmatrix}0&-1\\1&0\end{bmatrix}
  = \begin{bmatrix}0&-1\\-1&0\end{bmatrix},
\]
which is the reflection in the line $x_2=-x_1$. In the other order,
$[R_{\pi/2}][F]=\begin{bmatrix}0&1\\1&0\end{bmatrix}$, the reflection in $x_2=x_1$; order matters.
Also $[R_\alpha][R_\beta]=[R_{\alpha+\beta}]$, which (multiplying out) re-proves the angle-sum formulas.
\end{example}

\section{Range and kernel}

\begin{definition}
For a linear mapping $L:\R^n\to\R^m$,
\[
  \Range(L)=\{L(\vec x):\vec x\in\R^n\}\subseteq\R^m, \qquad
  \Ker(L)=\{\vec x\in\R^n : L(\vec x)=\vec 0\}\subseteq\R^n.
\]
\end{definition}

\begin{theorem}
$\Range(L)$ is a subspace of $\R^m$ and $\Ker(L)$ is a subspace of $\R^n$. If $A=[L]$,
$\Range(L)=\Span\{\vec a_1,\dots,\vec a_n\}$ (the column space of $A$) and $\Ker(L)$ is the solution
space of $A\vec x=\vec 0$.
\end{theorem}
\begin{proof}
$L(\vec0)=\vec0$ puts $\vec 0$ in both. If $L(\vec x)=L(\vec y)=\vec 0$ then
$L(\vec x+c\vec y)=\vec 0+c\vec 0=\vec 0$. If $\vec u=L(\vec x),\vec v=L(\vec y)$ then
$\vec u+c\vec v=L(\vec x+c\vec y)\in\Range(L)$. The descriptions follow from
$L(\vec x)=A\vec x = x_1\vec a_1+\dots+x_n\vec a_n$.
\end{proof}

\begin{example}
$L(\vec x)=A\vec x$ with $A=\begin{bmatrix}1&1&2\\1&2&3\end{bmatrix}$.
\[
\begin{bmatrix}1&1&2\\1&2&3\end{bmatrix}\xrightarrow{R_2-R_1}\begin{bmatrix}1&1&2\\0&1&1\end{bmatrix}
\xrightarrow{R_1-R_2}\begin{bmatrix}1&0&1\\0&1&1\end{bmatrix}.
\]
Kernel: $x_3=t$, $x_1=x_2=-t$, so $\Ker(L)=\Span\{(-1,-1,1)\}$. Since $\rank A=2=m$, $A\vec x=\vec b$
is consistent for all $\vec b$, so $\Range(L)=\R^2$.
\end{example}

\begin{example}
$L(x_1,x_2)=(x_1+x_2,\ 2x_1+2x_2,\ x_1+x_2)$ has $[L]=\begin{bmatrix}1&1\\2&2\\1&1\end{bmatrix}$, so
$\Range(L)=\Span\{(1,2,1)\}$ (a line in $\R^3$) and $\Ker(L)=\Span\{(1,-1)\}$.
\end{example}

\begin{important}[Geometric kernels and ranges]
For $\proj_{\vec d}$ on $\R^2$: range = the line $\Span\{\vec d\}$, kernel = the line orthogonal to
$\vec d$. Rotations and reflections have kernel $\{\vec0\}$ and range all of $\R^2$.
\end{important}

\begin{theorem}
$L$ is one-to-one iff $\Ker(L)=\{\vec 0\}$. $L:\R^n\to\R^m$ is onto iff $\rank[L]=m$.
\end{theorem}
\begin{proof}[Proof of the first claim]
If $L$ is one-to-one and $L(\vec x)=\vec 0=L(\vec 0)$ then $\vec x=\vec 0$. Conversely if the kernel
is trivial and $L(\vec x)=L(\vec y)$, then $L(\vec x-\vec y)=\vec 0$ so $\vec x=\vec y$.
\end{proof}

\section*{Summary}
\begin{itemize}
  \item Linear: $L(b\vec x+c\vec y)=bL(\vec x)+cL(\vec y)$; then $L(\vec0)=\vec0$.
  \item Every linear $L:\R^n\to\R^m$ is $\vec x\mapsto[L]\vec x$ with $[L]=[L(\vec e_1)\cdots L(\vec e_n)]$.
  \item $[R_\theta]$, $[\proj_{\vec d}]=\vec d\vec d^T/\norm{\vec d}^2$, $[\refl_{\vec d}]=2[\proj_{\vec d}]-I$.
  \item $[L\circ M]=[L][M]$: composition is matrix multiplication (order matters).
  \item $\Range(L)$ = span of the columns of $[L]$; $\Ker(L)$ = solutions of $[L]\vec x=\vec0$; both
        are subspaces. One-to-one iff trivial kernel.
\end{itemize}
